\section{Challenges and Motivation}
\label{sec:goal}
None of the three timelines can be reconstructed from its own vantage point: the OS sees threads but not what they run, an executor sees callbacks but not where they run, and the GPU sees work but not whose it is. The two kinds of information Section~\ref{sec:intro} asks for are therefore two challenges: the application architecture has to be complete, and the timelines have to be mapped onto it. Existing tools fall short on both: they recover the architecture partially, and none records the GPU.

\noindent\textbf{A complete application architecture.} The architecture is the core of every timeline and necessary to understand functional and timing behaviour: which components the software consists of, which functions are called when those components are triggered, and which channels they use to communicate. Knowing it completely means knowing four things. Which executor serves which nodes, and how each executor is configured: single- or multi-threaded, with how many threads, and with which callback groups, which decide which callbacks are allowed to run in parallel. The kind of every dispatchable unit of work---timer, subscription, service, client or waitable---since the executor treats them differently. Every communication pattern between callbacks: producer--consumer over a topic, request--response over a service, sample-and-hold, where a callback reads the latest value another one stored, and join, where a callback fires only once a set of inputs is complete. And the mechanism behind each: a serialised message through the middleware, a pointer handed over inside the process, or state that one callback stores and another later reads, which never becomes a message and can only be inferred.

Neither source code nor current tools capture this architecture. Source code does not contain what is decided during execution: which objects a given run creates, whether a delivery stays inside the process, and which callback publishes on which topic---a publisher belongs to the node, not to a callback, and any callback of the node may publish through it. Since composable nodes are loaded into containers at run time and handles are reused as objects come and go, even the object list has to be observed rather than read. Tools that trace the running system recover parts. ROS-Llama~\cite{blass-2021-rosllama-paper}, Abaza et al.~\cite{abaza-2024-ros2modelsynthesis} and LiME$^{\text{IT}}$~\cite{brandenburg-2026-limeit} recover callbacks and their execution times but assume a single-threaded executor; CARET and its extension~\cite{kuboichi-2022-caret,peng-2022-caretcomp} record nodes, callbacks, topics and executors but not the thread a callback runs on; and ros2\_tracing, on which they build, provides no tracepoints for waitables---the fifth kind the executor dispatches---or for deliveries inside a process, which never reach the client library where its hooks live. Executor configuration is recorded by none of them.

\noindent\textbf{Timelines mapped onto the architecture.} Where and when each piece of the architecture ran is decided on the three timelines of Fig.~\ref{fig:timelines}, and each of them has to be tied to the callback execution it belongs to. On the OS timeline the tie is the thread: under a multi-threaded executor the thread a callback runs on is decided at run time, and without the exact thread of every execution neither what the OS did to it---kept it on a core, preempted it, or let it block---nor anything else the thread did meanwhile can be attributed. The GPU is a second domain: it is traced by different tools on a different clock, what happens on it is invisible from the CPU, and it is a scheduling domain of its own, with compute and copy engines, each with its own queue, and a scheduler that orders the submitted work by stream and priority, unseen by both the OS and the executor. A callback that offloads work is then no longer atomic: what ROS~2 and the OS see as one execution becomes an alternating sequence of CPU segments and accelerator segments~\cite{enright-2024-paam}, and what a CPU-side trace records as its duration mixes computation with waiting for the GPU. Its segments are not fixed either: the GPU work a callback submits branches from invocation to invocation, so the sequence of segments---and hence the timing---differs between invocations of the same callback. And the only link between the two domains is again the thread the GPU run-time was called from. No existing extractor records the accelerator at all, which leaves the time a callback waits for its GPU work indistinguishable from computation.

The mapping is moreover not the same throughout a run. Some communications and executions occur only while the application initialises, some only while it shuts down, and the periodic operation in between is what matters from a real-time point of view; some registered callbacks never run at all under a given input. The architecture does not carry this distinction; it has to be recovered from what actually executed, and the phases separated before anything in the steady state is measured.

The two together are what end-to-end timing needs. A task chain, from a sensor reading to an actuation, is a path through the architecture that the timelines show actually ran, and it crosses executors, nodes and processes and both domains. Nor does a chain announce itself: in a deployment of hundreds of nodes the graph is entangled, and existing tools leave the chain to the user---CARET measures a path the user names, and the extractors above produce a callback graph without the phases or the GPU segments a chain runs through. What is needed is one model that holds the architecture and the three timelines with the relations of Fig.~\ref{fig:timelines} that tie them together, recorded on one clock and separated by phase, so that any timeline can be reconstructed alone and all of them in combination. Such a model serves more than analysis: it lets a traced run be revisited at any granularity, and it lets the same application be re-examined under a platform the trace did not run on---another core count, scheduler, executor configuration or GPU. All of it has to be obtained without modifying the application, and a model that is observed rather than assumed is only as good as its coverage, so the extraction itself must be checked. Section~\ref{sec:method} describes how FISH builds it.
